% --------------------------------------------------
% 6. ROBUSTNESS CHECKS
% --------------------------------------------------
\section{Robustness Checks}
\label{sec:robustness}

\subsection{Sub-sample Analysis: GC Contenders and Non-GC Riders}

A natural concern with the pooled model is that the determinants of abandonment may differ between GC contenders and the rest of the field. To investigate this, separate Cox PH models are estimated for GC contenders ($n = 1{,}395$) and non-GC riders ($n = 7{,}650$). The full regression tables are reported in Tables~\ref{tab:cox_gc} and~\ref{tab:cox_nonGC} in the Appendix; here the key findings are summarised.

The previous DNF rate effect is substantially larger for GC contenders (HR $= 2.280$, $p = 0.001$) than for non-GC riders (HR $= 1.472$, $p < 0.001$), and the confidence intervals barely overlap, suggesting meaningful heterogeneity. This is consistent with a higher-stakes signalling environment for GC contenders: a rider whose past abandonment record has already signalled fragility faces compounding reputational and physical costs. In tournament theory terms \parencite{rosen1981superstars, lazear1981rank}, the amplified effect may also reflect the sharper incentive cliff that GC contenders face once they fall out of contention.

The Tour de France protective effect is significant only for non-GC riders (HR $= 0.866$, $p = 0.005$), not for GC contenders (HR $= 0.863$, $p = 0.177$). Although the point estimates are nearly identical, the GC estimate loses precision due to the smaller sub-sample. The pattern suggests that event-value incentives operate primarily through domestiques and support riders, for whom finishing the Tour carries direct benefits in team bonuses, sponsor visibility, and reputational value. GC contenders face a more binary decision: continue if competitive, withdraw if not, regardless of which Grand Tour they are racing.

Economically, this heterogeneity has implications for how we interpret the pooled model. The pooled GC contender coefficient (HR $= 1.156$) is a weighted average of two distinct behavioural regimes: one in which GC contenders are driven by tournament incentives and physical exposure, and one in which non-GC riders respond to event-value gradients. Pooling conceals the fact that the economic forces operating on these two groups are qualitatively different, not just quantitatively. This suggests that future work modelling abandonment should allow for interaction effects between rider role and race-level incentives, rather than assuming a single set of coefficients applies to all riders. The sub-sample results also reinforce the practical value of the analysis for team management: a team selecting riders for a Grand Tour roster should weight a GC candidate's abandonment history more heavily than a domestique's, because the predictive power of past DNF behaviour is substantially stronger for the former group.
